\section{Background and Lineage}
\label{sec:background}

The transition from static algorithms to autonomous software engineering agents is underpinned by a century-long chain of intellectual developments in computation, information, learning, scale, and distributed systems. Understanding this lineage is essential to framing agentic software engineering not as an ad-hoc scripting trend, but as a formal systems execution domain.

\subsection{Formalizing the Machine and the Message}
The mathematical foundation of computing begins with Turing's formalization of computation in 1936~\cite{turing1936computable}. By introducing the Turing Machine, Alan Turing defined the limits of what can be computed algorithmically and established that a universal machine could simulate any other process given a discrete set of rules and an infinite tape. 

Twelve years later, Shannon defined what flows through these machines~\cite{shannon1948mathematical}. In \emph{A Mathematical Theory of Communication}, Claude Shannon stripped meaning from language, formalizing information as a measurable quantity of entropy and establishing the bit as the core unit of channel capacity. This mathematical framing of entropy and prediction under uncertainty directly underpins the loss functions and token-generation architectures used in modern language models.

\subsection{The Connectionist Evolution and Distributed Ordering}
Early attempts to make machines learn from their environments began with Rosenblatt's Perceptron in 1958~\cite{rosenblatt1958perceptron}, which showed that machines could adjust numerical weights to classify input patterns. However, Minsky and Papert's critique in 1969~\cite{minsky1969perceptrons} exposed the geometric limitations of single-layer networks, notably their inability to learn non-linearly separable functions like XOR. This critique initiated the first AI winter, which was resolved seventeen years later by Rumelhart, Hinton, and Williams~\cite{rumelhart1986learning} who formalized backpropagation, showing that stacking layers of perceptrons and using the chain rule from calculus allowed multi-layered networks to learn complex internal representations.

Concurrently, the systems substrate of computing was expanding from single processors to distributed networks. Lamport's seminal 1978 paper~\cite{lamport1978time} established that without a shared physical clock, events in a distributed system must be ordered based on causality (the "happens-before" relation) rather than wall-clock time. This conceptual breakthrough made large-scale distributed databases, cloud storage, and eventually parallel GPU training clusters possible, preventing concurrent machines from dissolving into state inconsistency.

\subsection{Scale, Attention, and Emergent Intelligence}
In 1998, Brin and Page demonstrated that relevance could emerge from planetary-scale graph structure through the PageRank algorithm~\cite{brin1998anatomy}. This proved that indexing and ranking massive piles of human text and links could yield high-quality search interfaces. By 2012, the ImageNet breakthrough (AlexNet) by Krizhevsky et al.~\cite{krizhevsky2012imagenet} proved that connectionist backpropagation, when combined with massive datasets and consumer-grade graphics processing units (GPUs), dramatically outperformed hand-engineered features in computer vision.

The modern era of LLMs was ignited by Vaswani et al. in 2017~\cite{vaswani2017attention} with the Transformer architecture. By discarding sequential recurrent networks in favor of parallelizable attention mechanisms, the Transformer allowed models to compute associations between all tokens in a context window simultaneously. Finally, in 2020, Brown et al. showed in GPT-3~\cite{brown2020language} that scaling these transformers to hundreds of billions of parameters transforms language predictors into general-purpose few-shot learners. GPT-3 proved that scale turns predictions into programmable interfaces.

\subsection{The Agentic Transition and the Accountability Gap}
While GPT-3 and subsequent models excel at static text completion, software engineering requires active intervention. This transition from static representation to dynamic action was foreshadowed by Mnih et al. in Deep Q-Networks (DQN)~\cite{mnih2015human}, where neural networks were integrated into closed-loop reinforcement learning environments—mapping perception directly to actions and environment feedback.

Today, LLM agents combine the semantic reasoning of transformers with the action loops of reinforcement learning, executing commands via shell tools over mutable code repositories. Yet, as computation has transitioned from Turing's abstract tape, through Shannon's bits, Lamport's clocks, and Transformer attention, we have reached a critical systems gap: \textbf{we have executable intent, but we lack custody and proof}. This paper bridges this gap, developing the control plane primitives needed to make autonomous agent runs safe and accountable.
